\documentclass{article}
\usepackage{amsmath,amssymb,amsthm}
\usepackage{algorithm}
\usepackage{algpseudocode}
\usepackage{hyperref}
\usepackage{enumitem}
\usepackage{geometry}
\geometry{margin=1in}
\newtheorem{theorem}{Theorem}
\newtheorem{lemma}{Lemma}
\newtheorem{assumption}{Assumption}
\newcommand{\prox}{\operatorname{prox}}
\newcommand{\R}{\mathbb{R}}
\newcommand{\E}{\mathbb{E}}

\begin{document}

\section*{Proximal Gradient Method (PGM)}

\section*{Mathematical Formulation}
We consider the composite convex optimisation problem  

\[
\min_{x\in\R^{d}} \; F(x)\;\; \text{where}\;\; F(x):=f(x)+g(x).
\]

\begin{itemize}
\item $f:\R^{d}\rightarrow\R$ is convex, differentiable and has an $L$‑Lipschitz continuous gradient:
\[
\|\nabla f(u)-\nabla f(v)\|\le L\|u-v\|,\qquad\forall u,v\in\R^{d}.
\]
\item $g:\R^{d}\rightarrow\R\cup\{+\infty\}$ is proper, closed and convex (possibly non‑smooth).
\item For a stepsize $\eta>0$, the proximal operator of $g$ is
\[
\prox_{\eta g}(z):=\arg\min_{x\in\R^{d}}\Big\{ g(x)+\frac{1}{2\eta}\|x-z\|^{2}\Big\}.
\]
\end{itemize}

The \textbf{proximal gradient update} reads  

\[
\boxed{
x_{k+1}= \prox_{\eta g}\bigl(x_{k}-\eta\nabla f(x_{k})\bigr)},\qquad k=0,1,\dots
\tag{1}
\]

The stepsize is chosen in the interval  

\[
\eta\in\bigl(0,\tfrac{1}{L}\bigr].
\tag{2}
\]

\section*{Pseudocode}
\begin{algorithm}[H]
\caption{Proximal Gradient Method}
\begin{algorithmic}[1]
\Require Initial point $x_{0}\in\R^{d}$, stepsize $\eta\in(0,1/L]$, 
maximum iterations $K$, tolerance $\epsilon>0$.
\For{$k=0$ \textbf{to} $K-1$}
    \State $g_{k}\gets \nabla f(x_{k})$ \Comment{gradient of the smooth part}
    \State $z_{k}\gets x_{k}-\eta\,g_{k}$ \Comment{gradient step}
    \State $x_{k+1}\gets \prox_{\eta g}(z_{k})$ \Comment{proximal map}
    \If{$\|x_{k+1}-x_{k}\|\le\epsilon$}
        \State \textbf{break}
    \EndIf
\EndFor
\State \Return $x_{k+1}$
\end{algorithmic}
\end{algorithm}

\section*{Dry Run Example}
Let  

\[
f(w)=(w-3)^{2},\qquad g(w)=0\;\;(\text{so }\prox_{\eta g}= \text{id}),
\]
and initialise $w_{0}=0$, stepsize $\eta=0.1$.
The optimal solution is $w^{*}=3$.

\begin{center}
\begin{tabular}{c|c|c|c}
Iteration $k$ & $\nabla f(w_{k})$ & $w_{k+1}=w_{k}-\eta\nabla f(w_{k})$ & $F(w_{k+1})$\\ \hline
0 & $2(w_{0}-3) = -6$ & $0-0.1(-6)=0.6$ & $(0.6-3)^{2}=5.76$\\
1 & $2(0.6-3)=-4.8$ & $0.6-0.1(-4.8)=1.08$ & $(1.08-3)^{2}=3.6864$\\
2 & $2(1.08-3)=-3.84$ & $1.08-0.1(-3.84)=1.464$ & $(1.464-3)^{2}=2.3660$\\
\end{tabular}
\end{center}

The iterates move monotonically toward $3$ and the objective value decreases at each step.

\newpage
\section*{Assumptions}
\begin{assumption}[Smoothness of $f$]\label{as:lip}
$f$ is convex and its gradient is $L$‑Lipschitz (cf. (1) in the formulation).
\end{assumption}
\begin{assumption}[Convexity of $g$]\label{as:convexg}
$g$ is proper, closed and convex; its proximal operator $\prox_{\eta g}$ is well defined for all $\eta>0$.
\end{assumption}
\begin{assumption}[Stepsize]\label{as:stepsize}
The stepsize satisfies $0<\eta\le 1/L$.
\end{assumption}
All results below rely exclusively on Assumptions \ref{as:lip}–\ref{as:stepsize}.

\section*{Main Convergence Theorem}
\begin{theorem}[Function‑value convergence of PGM]\label{thm:main}
Let $\{x_{k}\}_{k\ge0}$ be generated by \eqref{1} with a stepsize $\eta\in(0,1/L]$.  
Denote $x^{*}\in\arg\min F$ and assume such a minimiser exists. Then for every $k\ge1$  

\[
F(x_{k})-F(x^{*})\;\le\;\frac{\|x_{0}-x^{*}\|^{2}}{2\eta\,k}.
\tag{3}
\]

Consequently $F(x_{k})\to F(x^{*})$ at the rate $O(1/k)$.
\end{theorem}

\section*{Theoretical Properties}
\begin{itemize}
\item \textbf{Stability.} Because the proximal operator is non‑expansive, the iterates cannot diverge as long as the gradient step is bounded (ensured by $\eta\le 1/L$).
\item \textbf{Hyperparameter sensitivity.} The convergence constant in \eqref{3} is inversely proportional to $\eta$; a larger admissible stepsize (up to $1/L$) yields a tighter bound.
\item \textbf{Relation to baselines.}
    \begin{enumerate}[label=(\alph*)]
        \item If $g\equiv0$, $\prox_{\eta g}=\text{id}$ and \eqref{1} reduces to classical gradient descent, and \eqref{3} coincides with the standard $O(1/k)$ bound.
        \item If $f\equiv0$, the method becomes the proximal point algorithm, which also satisfies \eqref{3} with $L=\infty$ (any $\eta>0$).
    \end{enumerate}
\end{itemize}

\section*{Discussion}
The theorem shows that the proximal gradient method inherits the optimal $O(1/k)$ rate for general convex composite problems, matching the best known rates without requiring strong convexity. The proof hinges on two elementary properties: (i) the descent lemma for $L$‑smooth functions and (ii) the non‑expansiveness of the proximal map. No stochasticity or higher‑order information is needed.

\newpage
\section*{Preliminaries (Definitions and Lemmas)}
\begin{lemma}[Descent Lemma]\label{lem:descent}
If $\nabla f$ is $L$‑Lipschitz, then for all $u,v\in\R^{d}$  

\[
f(v)\le f(u)+\langle\nabla f(u),v-u\rangle+\frac{L}{2}\|v-u\|^{2}.
\tag{4}
\]
\end{lemma}
\begin{proof}
From the $L$‑Lipschitz property and convexity of $f$; a standard integration argument yields \eqref{4}. \hfill $\square$
\end{proof}

\begin{lemma}[Non‑expansiveness of the proximal operator]\label{lem:nonexp}
For any $\eta>0$ and any $u,v\in\R^{d}$  

\[
\big\|\prox_{\eta g}(u)-\prox_{\eta g}(v)\big\|\le\|u-v\|.
\tag{5}
\]
\end{lemma}
\begin{proof}
The proximal map is the resolvent of the monotone sub‑differential $\partial g$; monotonicity implies \eqref{5}. \hfill $\square$
\end{proof}

\section*{Main Proof (Step‑by‑Step Derivation)}
We prove Theorem \ref{thm:main}. Throughout the proof we fix an arbitrary optimal point $x^{*}\in\arg\min F$.

\begin{enumerate}
\item[Step 1.]  Write the update \eqref{1} as a two‑stage operation:
\[
z_{k}:=x_{k}-\eta\nabla f(x_{k}),\qquad
x_{k+1}= \prox_{\eta g}(z_{k}). \tag{6}
\]

\item[Step 2.]  By the definition of the proximal operator (optimality condition) we have
\[
0\in \partial g(x_{k+1})+\frac{1}{\eta}(x_{k+1}-z_{k}),
\]
hence
\[
\frac{1}{\eta}(z_{k}-x_{k+1})\in\partial g(x_{k+1}). \tag{7}
\]

\item[Step 3.]  Use convexity of $g$ with point $x^{*}$:
\[
g(x^{*})\ge g(x_{k+1})+\Big\langle \frac{1}{\eta}(z_{k}-x_{k+1}),\,x^{*}-x_{k+1}\Big\rangle.
\tag{8}
\]

\item[Step 4.]  Rearrange \eqref{8} to obtain an inequality for $g(x_{k+1})-g(x^{*})$:
\[
g(x_{k+1})-g(x^{*})\le
\frac{1}{\eta}\langle z_{k}-x_{k+1},\,x_{k+1}-x^{*}\rangle.
\tag{9}
\]

\item[Step 5.]  Expand the inner product in \eqref{9} using $z_{k}=x_{k}-\eta\nabla f(x_{k})$:
\[
\langle z_{k}-x_{k+1},\,x_{k+1}-x^{*}\rangle
= \langle x_{k}-x_{k+1},\,x_{k+1}-x^{*}\rangle
-\eta\langle\nabla f(x_{k}),\,x_{k+1}-x^{*}\rangle.
\tag{10}
\]

\item[Step 6.]  Apply the identity $2\langle a,b\rangle=\|a\|^{2}+\|b\|^{2}-\|a-b\|^{2}$ to the first term of \eqref{10}:
\[
\langle x_{k}-x_{k+1},x_{k+1}-x^{*}\rangle
=\frac{1}{2}\Big(\|x_{k}-x^{*}\|^{2}-\|x_{k+1}-x^{*}\|^{2}-\|x_{k}-x_{k+1}\|^{2}\Big).
\tag{11}
\]

\item[Step 7.]  Substitute \eqref{11} into \eqref{10} and then into \eqref{9}:
\[
\begin{aligned}
g(x_{k+1})-g(x^{*})
\le &\;\frac{1}{2\eta}\Big(\|x_{k}-x^{*}\|^{2}-\|x_{k+1}-x^{*}\|^{2}
-\|x_{k}-x_{k+1}\|^{2}\Big)\\
&\;-\langle\nabla f(x_{k}),\,x_{k+1}-x^{*}\rangle.
\end{aligned}
\tag{12}
\]

\item[Step 8.]  Add $f(x_{k+1})-f(x^{*})$ to both sides of \eqref{12}. Using the descent lemma \eqref{4} with $u=x_{k}$, $v=x_{k+1}$ we have
\[
f(x_{k+1})\le f(x_{k})+\langle\nabla f(x_{k}),x_{k+1}-x_{k}\rangle
+\frac{L}{2}\|x_{k+1}-x_{k}\|^{2}.
\tag{13}
\]

\item[Step 9.]  Combine \eqref{12} and \eqref{13} to bound the total objective difference:
\[
\begin{aligned}
F(x_{k+1})-F(x^{*})
=&\;[f(x_{k+1})-f(x^{*})]+[g(x_{k+1})-g(x^{*})]\\[4pt]
\le&\;f(x_{k})-f(x^{*})+\langle\nabla f(x_{k}),x_{k+1}-x_{k}\rangle
+\frac{L}{2}\|x_{k+1}-x_{k}\|^{2}\\
&\;+\frac{1}{2\eta}\Big(\|x_{k}-x^{*}\|^{2}-\|x_{k+1}-x^{*}\|^{2}
-\|x_{k}-x_{k+1}\|^{2}\Big)\\
&\;-\langle\nabla f(x_{k}),x_{k+1}-x^{*}\rangle .
\end{aligned}
\tag{14}
\]

\item[Step 10.]  Observe that the gradient inner products cancel:
\[
\langle\nabla f(x_{k}),x_{k+1}-x_{k}\rangle
-\langle\nabla f(x_{k}),x_{k+1}-x^{*}\rangle
= -\langle\nabla f(x_{k}),x_{k}-x^{*}\rangle .
\tag{15}
\]

\item[Step 11.]  Drop the non‑positive term $-\langle\nabla f(x_{k}),x_{k}-x^{*}\rangle$ (it can be omitted because it only makes the RHS larger). After simplification we obtain
\[
F(x_{k+1})-F(x^{*})
\le \frac{1}{2\eta}\bigl(\|x_{k}-x^{*}\|^{2}-\|x_{k+1}-x^{*}\|^{2}\bigr)
+\Big(\frac{L}{2}-\frac{1}{2\eta}\Big)\|x_{k+1}-x_{k}\|^{2}.
\tag{16}
\]

\item[Step 12.]  By Assumption \ref{as:stepsize}, $\eta\le 1/L$, hence $\frac{L}{2}-\frac{1}{2\eta}\le0$. Therefore the second term in \eqref{16} is non‑positive and can be discarded, yielding the fundamental inequality
\[
F(x_{k+1})-F(x^{*})
\le \frac{1}{2\eta}\bigl(\|x_{k}-x^{*}\|^{2}-\|x_{k+1}-x^{*}\|^{2}\bigr).
\tag{17}
\]

\item[Step 13.]  Sum \eqref{17} over $k=0,\dots ,K-1$ (telescoping):
\[
\sum_{k=0}^{K-1}\bigl(F(x_{k+1})-F(x^{*})\bigr)
\le \frac{1}{2\eta}\bigl(\|x_{0}-x^{*}\|^{2}-\|x_{K}-x^{*}\|^{2}\bigr)
\le \frac{\|x_{0}-x^{*}\|^{2}}{2\eta}.
\tag{18}
\]

\item[Step 14.]  Since $F(x_{k})$ is non‑increasing (from \eqref{17}), each term in the sum is bounded above by the average:
\[
K\bigl(F(x_{K})-F(x^{*})\bigr)
\le \sum_{k=0}^{K-1}\bigl(F(x_{k+1})-F(x^{*})\bigr).
\tag{19}
\]

\item[Step 15.]  Combine \eqref{18} and \eqref{19} and replace $K$ by any $k\ge1$:
\[
F(x_{k})-F(x^{*})\le\frac{\|x_{0}-x^{*}\|^{2}}{2\eta\,k},
\qquad k\ge1.
\tag{20}
\]

\item[Step 16.]  Inequality \eqref{20} is exactly the claim \eqref{3} of Theorem \ref{thm:main}. \hfill $\blacksquare$
\end{enumerate}

\section*{Rate Derivation (With Constants)}
From \eqref{20} we read off the explicit convergence constant:
\[
C:=\frac{\|x_{0}-x^{*}\|^{2}}{2\eta},
\qquad\text{so that}\qquad
F(x_{k})-F(x^{*})\le \frac{C}{k}.
\]
If the maximal admissible stepsize $\eta=1/L$ is used, then $C=\frac{L}{2}\|x_{0}-x^{*}\|^{2}$.

\section*{Special Cases}
\begin{itemize}
\item \textbf{Pure gradient descent} ($g\equiv0$): $\prox_{\eta g}= \text{id}$. The same proof yields \eqref{3} with $g\equiv0$.
\item \textbf{Proximal point algorithm} ($f\equiv0$): the update reduces to $x_{k+1}= \prox_{\eta g}(x_{k})$. Inequality \eqref{17} becomes $F(x_{k+1})-F(x^{*})\le\frac{1}{2\eta}(\|x_{k}-x^{*}\|^{2}-\|x_{k+1}-x^{*}\|^{2})$, giving the same $O(1/k)$ bound without any smoothness assumption.
\item \textbf{Strongly convex $F$}: if $F$ is $\mu$‑strongly convex, a standard refinement of \eqref{17} leads to a linear rate $F(x_{k})-F(x^{*})\le (1-\eta\mu)^{k}\frac{\|x_{0}-x^{*}\|^{2}}{2\eta}$.
\end{itemize}

\end{document}